\chapter{Forecasting Monthly Sales}
\label{ch:forecasting}

This is the most theory-heavy chapter of the report: forecasting a short
(48-point) monthly series correctly requires care about stationarity,
model identification, and — critically — an evaluation protocol that
respects time order, none of which arise in the i.i.d.\ supervised setting
of \cref{ch:models}.

\section{Theory}

\subsection{Stationarity}
\begin{definition}[Weak stationarity]
A process $\{y_t\}$ is \emph{weakly (covariance) stationary} if
$\mathbb{E}[y_t] = \mu$ is constant in $t$, $\mathrm{Var}(y_t)=\sigma^2$
is constant in $t$, and $\mathrm{Cov}(y_t, y_{t+h})$ depends on the lag
$h$ but not on $t$ itself.
\end{definition}
Classical linear time-series models (AR, MA, ARMA) are defined for
stationary processes; a series with a trend or a growing seasonal
amplitude violates the definition and must be transformed — typically by
\emph{differencing} — before such a model applies. The backshift (lag)
operator $B$, $By_t = y_{t-1}$, lets a $d$-th order difference be written
compactly as $\Delta^d y_t = (1-B)^d y_t$; a \emph{seasonal} difference of
period $s$ is $\Delta_s y_t = (1-B^s) y_t = y_t - y_{t-s}$.

\subsection{The augmented Dickey--Fuller test}
The ADF test formalises "does this series need differencing" as a
hypothesis test on the regression
\begin{equation}
  \Delta y_t = \alpha + \beta t + \gamma y_{t-1}
             + \sum_{i=1}^{k}\delta_i \Delta y_{t-i} + \varepsilon_t,
  \label{eq:adf}
\end{equation}
with null hypothesis $H_0: \gamma = 0$ (a unit root — the series is a
random walk, non-stationary) against $H_1: \gamma < 0$ (stationary). The
test statistic is the ordinary $t$-statistic for $\hat\gamma$, but its
null distribution is \emph{not} the usual Student-$t$ (it is skewed left
by the non-standard behaviour of a random walk's sample variance), so
Dickey and Fuller's own tabulated critical values (or the software's
$p$-value approximation) must be used rather than a normal-theory table.

\subsection{Autocorrelation and the Box--Jenkins identification}
The autocorrelation function $\rho(h) = \mathrm{Corr}(y_t, y_{t+h})$ and
partial autocorrelation function $\phi_{hh}$ (the correlation between
$y_t$ and $y_{t+h}$ after regressing out $y_{t+1},\dots,y_{t+h-1}$) are
the classical diagnostic pair for identifying ARMA order: an AR($p$)
process has a PACF that cuts off after lag $p$ and an ACF that decays
geometrically, while an MA($q$) process has the mirror-image behaviour.
In practice for a short, strongly seasonal series like ours, ACF/PACF are
used qualitatively (\cref{fig:fc-acf-pacf}) to corroborate, not to
mechanically dictate, the model order.

\subsection{Exponential smoothing (Holt--Winters)}
Simple exponential smoothing forecasts the next value as a
geometrically-weighted average of the past, $\hat y_{t+1} = \alpha y_t +
(1-\alpha)\hat y_t$. \textbf{Holt--Winters} extends this with a trend
state $b_t$ and a seasonal state $s_t$ (period $L=12$). The additive
form used here is
\begin{align}
  \ell_t &= \alpha (y_t - s_{t-L}) + (1-\alpha)(\ell_{t-1}+b_{t-1}), \label{eq:hw-level}\\
  b_t &= \beta (\ell_t - \ell_{t-1}) + (1-\beta) b_{t-1}, \label{eq:hw-trend}\\
  s_t &= \gamma (y_t - \ell_t) + (1-\gamma) s_{t-L}, \label{eq:hw-season}\\
  \hat y_{t+h} &= \ell_t + h\,b_t + s_{t-L+((h-1)\bmod L)+1}, \label{eq:hw-forecast}
\end{align}
with smoothing parameters $\alpha,\beta,\gamma \in [0,1]$ chosen (by the
fitting routine) to minimise the in-sample sum of squared one-step
forecast errors. Level $\ell_t$ tracks the de-seasonalised local mean,
$b_t$ the local trend, and $s_t$ a smoothly-updating length-$L$ seasonal
profile; the $h$-step forecast \cref{eq:hw-forecast} linearly extrapolates
the trend and re-applies the appropriate seasonal index. This additive
form is appropriate when the seasonal swing has roughly constant absolute
size rather than growing proportionally with the level (a multiplicative
alternative replaces the $+s$ terms with $\times s$ terms), consistent
with \cref{fig:stl} where the seasonal amplitude does not visibly balloon
with the trend.

\subsection{SARIMA}
A seasonal ARIMA model, $\mathrm{SARIMA}(p,d,q)(P,D,Q)_s$, generalises
ARMA by combining non-seasonal and seasonal autoregressive/moving-average
polynomials applied to a series differenced both regularly ($d$ times)
and seasonally ($D$ times, period $s$):
\begin{equation}
  \phi_p(B)\,\Phi_P(B^s)\, (1-B)^d (1-B^s)^D y_t
    = \theta_q(B)\,\Theta_Q(B^s)\,\varepsilon_t,
  \label{eq:sarima}
\end{equation}
where $\phi_p(B) = 1-\phi_1 B - \cdots - \phi_p B^p$ (and analogously
$\Phi_P$, $\theta_q$, $\Theta_Q$ for the seasonal AR, non-seasonal MA and
seasonal MA polynomials respectively) and $\varepsilon_t$ is white noise.
We fit $\mathrm{SARIMA}(1,1,1)(1,1,0,12)$: one non-seasonal AR and MA
term, one seasonal AR term, one regular and one seasonal (period-12)
difference, and no seasonal MA term — a compact specification chosen to
be identifiable on 48 points while still capturing both the trend
($d=1$) and the year-over-year seasonal shape ($D=1$, period 12) directly,
rather than relying on Holt--Winters' exponential-smoothing-style
seasonal update. Coefficients are estimated by (quasi-)maximum likelihood
via the Kalman filter (as implemented in \texttt{statsmodels}'s
state-space \texttt{SARIMAX}), which also yields exact prediction
variances and hence analytic confidence intervals for the forecast.

\subsection{Forecast accuracy metrics}
For actuals $y_1,\dots,y_h$ and forecasts $\hat y_1,\dots,\hat y_h$ over a
horizon of length $h$:
\begin{align}
  \mathrm{MAE} &= \frac1h\sum_{i=1}^h |y_i - \hat y_i|, &
  \mathrm{RMSE} &= \sqrt{\frac1h\sum_{i=1}^h (y_i-\hat y_i)^2}, \\
  \mathrm{MAPE} &= \frac{100\%}{h}\sum_{i=1}^h \left|\frac{y_i-\hat y_i}{y_i}\right|. &&
\end{align}
MAE is in the original (dollar) units and treats all errors linearly;
RMSE is also in dollars but penalises large errors quadratically, so it
is more sensitive to occasional big misses (e.g. an under-forecast
December); MAPE is scale-free and hence comparable across series of
different magnitude (used here to compare the whole-store series against
per-category series of very different size), at the cost of being
undefined/unstable when $y_i$ is near zero and of asymmetrically
penalising over- vs.\ under-forecasts.

\subsection{Rolling-origin (walk-forward) backtesting}
A single train/test split, as used for the i.i.d.\ models of
\cref{ch:models}, is a poor evaluation protocol for a 48-point time
series: it tests the model at exactly one forecast origin, so its score
has high variance and cannot be trusted to represent typical future
performance. \textbf{Rolling-origin} (a.k.a.\ walk-forward, or expanding-
window) backtesting instead repeats the fit/forecast cycle at multiple
origins $o_1 < o_2 < \dots < o_K$: at each $o_k$, the model is trained
only on $y_1,\dots,y_{o_k}$ and scored against the next $h$ true values
$y_{o_k+1},\dots,y_{o_k+h}$, and the reported metric is the average
across the $K$ folds. This never trains on the future relative to any
fold's own test window and is the time-series analogue of $K$-fold
cross-validation, with the essential difference that folds must respect
temporal order rather than being drawn at random.

\subsection{Residual whiteness: the Ljung--Box test}
A well-specified model should leave residuals $\hat\varepsilon_t$ that
are approximately white noise (no remaining autocorrelation to exploit).
The Ljung--Box statistic
\begin{equation}
  Q_{LB} = n(n+2)\sum_{k=1}^{K}\frac{\hat\rho_k^2}{n-k}
  \label{eq:ljungbox}
\end{equation}
(with $\hat\rho_k$ the residual sample autocorrelation at lag $k$)
is asymptotically $\chi^2_K$ under the null of no autocorrelation up to
lag $K$; a large $p$-value (failure to reject) is the desired outcome,
indicating the model has not left exploitable structure behind.

\section{Method}
The monthly series $y_t = \sum_{\text{lines in month } t} \text{Sales}$
(\texttt{src.features.monthly\_sales}) has $T=48$ points, 2014-01 through
2017-12. We hold out the last 12 months as a single test window, fit
four forecasters on the code interface \texttt{fit(train)} /
\texttt{predict(h)} (\texttt{src/forecasting.py}) — a seasonal-naive
baseline, additive Holt--Winters, $\mathrm{SARIMA}(1,1,1)(1,1,0,12)$, and
a recursive random-forest regressor on lag/rolling-mean/seasonal features
(\texttt{LagRegressor}) — and separately run an 8-fold rolling-origin
backtest (\texttt{min\_train}=24, horizon=3, step=3). The winning point
forecaster is then refit on the full 48 months and projected 12 months
ahead; SARIMA additionally supplies an 80\% analytic confidence band via
its Kalman-filter forecast variance.

\section{Results}

\subsection{Series and stationarity}
\begin{figure}[H]
  \centering
  \includegraphics[width=0.95\linewidth]{04_series.png}
  \caption{Monthly sales, 2014--2017. Clear multi-year growth and a
  recurring within-year cycle.}
  \label{fig:fc-series}
\end{figure}
Table~\ref{tab:adf} reports the ADF test (\cref{eq:adf}) on the level
series, its first difference, and its seasonal (lag-12) difference.
\begin{table}[H]
  \centering
  \caption{ADF unit-root test, $H_0$: series has a unit root.}
  \label{tab:adf}
  \begin{tabular}{@{}lrr@{}}
    \toprule
    Series & Test statistic & $p$-value \\
    \midrule
    Level $y_t$                  & $-4.494$ & 0.0002 \\
    First difference $\Delta y_t$ & $-9.058$ & $<$0.0001 \\
    Seasonal difference $\Delta_{12} y_t$ & $-1.482$ & 0.5425 \\
    \bottomrule
  \end{tabular}
\end{table}
The level series already rejects a pure unit root, and first-differencing
rejects it far more decisively — but the \emph{seasonal-only} difference
does \textbf{not} (p = 0.54): removing only the year-over-year level
shift is not enough to leave a stationary residual, because the
within-year \emph{shape} of the cycle (not just its level) still carries
structure. This is exactly the case for using a model with an explicit,
separately-estimated seasonal component — Holt--Winters' $s_t$ state, or
SARIMA's combined regular-\emph{and}-seasonal differencing
(\cref{eq:sarima}, $d{=}1,D{=}1$) — rather than relying on a single
seasonal difference alone.
\begin{figure}[H]
  \centering
  \includegraphics[width=0.95\linewidth]{04_acf_pacf.png}
  \caption{ACF (left) and PACF (right) of $\Delta y_t$. A significant
  spike near lag 12 in both is consistent with the seasonal structure
  seen directly in \cref{fig:fc-series,fig:stl}.}
  \label{fig:fc-acf-pacf}
\end{figure}

\subsection{Hold-out comparison}
\begin{table}[H]
  \centering
  \caption{12-month hold-out accuracy (train on months 1--36, test on 37--48).}
  \label{tab:holdout}
  \begin{tabular}{@{}lrrr@{}}
    \toprule
    Model & MAE (\$) & RMSE (\$) & MAPE \\
    \midrule
    \textbf{Holt--Winters} & \textbf{11{,}456} & 12{,}544 & \textbf{22.6\%} \\
    SARIMA(1,1,1)(1,1,0,12) & 13{,}411 & 15{,}989 & 26.7\% \\
    Lag-regressor (RF)      & 13{,}484 & 16{,}791 & 24.1\% \\
    Seasonal-naive baseline & 15{,}468 & 18{,}986 & 24.5\% \\
    \bottomrule
  \end{tabular}
\end{table}
\begin{figure}[H]
  \centering
  \includegraphics[width=0.95\linewidth]{04_holdout_forecasts.png}
  \caption{12-month hold-out: all four models' forecasts against the
  realised sales path.}
  \label{fig:fc-holdout}
\end{figure}
All three fitted models beat the seasonal-naive baseline on MAE and RMSE;
Holt--Winters is the clear winner on this single split.

\subsection{Rolling-origin backtest}
\begin{table}[H]
  \centering
  \caption{8-fold rolling-origin backtest (\texttt{min\_train}=24,
  horizon=3, step=3), metrics averaged across folds.}
  \label{tab:backtest}
  \begin{tabular}{@{}lrrrr@{}}
    \toprule
    Model & MAE (\$) & RMSE (\$) & MAPE & Folds \\
    \midrule
    \textbf{SARIMA}         & \textbf{11{,}612} & 13{,}373 & 23.9\% & 8 \\
    Holt--Winters            & 12{,}314 & 14{,}104 & 25.1\% & 8 \\
    Lag-regressor (RF)       & 13{,}660 & 15{,}527 & 27.2\% & 8 \\
    Seasonal-naive baseline  & 13{,}994 & 16{,}207 & 24.9\% & 8 \\
    \bottomrule
  \end{tabular}
\end{table}
\begin{figure}[H]
  \centering
  \includegraphics[width=0.95\linewidth]{04_backtest_predictions.png}
  \caption{Rolling-origin backtest predictions (Holt--Winters and SARIMA)
  vs.\ actual, across all 8 folds.}
  \label{fig:fc-backtest}
\end{figure}
The ranking flips: with more forecast origins averaged, SARIMA edges out
Holt--Winters (\$11{,}612 vs.\ \$12{,}314 MAE), though the two remain
within about 6\% of each other and both clearly beat the naive baseline
and the lag-regressor. This is a useful illustration of exactly the
risk the theory section warns about: the single hold-out split
(\cref{tab:holdout}) and the 8-fold backtest (\cref{tab:backtest})
disagree on which of the top two models is "best", and the backtest,
having 8 origins rather than 1, is the more trustworthy estimate of
typical future accuracy.

\subsection{Final forecast}
Given the closeness of the two models and their complementary strengths
(Holt--Winters: best single-split point forecast; SARIMA: best
backtest average \emph{and} a principled analytic interval), the final
projection uses \textbf{Holt--Winters, refit on all 48 months, for the
point forecast} and \textbf{SARIMA for an 80\% confidence band}.
\begin{figure}[H]
  \centering
  \includegraphics[width=0.95\linewidth]{04_final_forecast.png}
  \caption{History plus the 12-month-ahead forecast: Holt--Winters point
  forecast and the SARIMA 80\% interval.}
  \label{fig:fc-final}
\end{figure}
The next 12 months are projected at \textbf{\$914{,}903}, against
\$733{,}215 realised in the last full year (2017) — a projected
\textbf{+24.8\%} year-on-year. Table~\ref{tab:cat-forecast} breaks this
down by product category (each fit with its own Holt--Winters model on
that category's monthly series).
\begin{table}[H]
  \centering
  \caption{Per-category 2018 projection (Holt--Winters).}
  \label{tab:cat-forecast}
  \begin{tabular}{@{}lrrr@{}}
    \toprule
    Category & 2017 actual (\$) & 2018 projected (\$) & Growth \\
    \midrule
    Office Supplies & 246{,}097 & 326{,}734 & \textbf{+32.8\%} \\
    Furniture       & 215{,}387 & 245{,}463 & +14.0\% \\
    Technology      & 271{,}731 & 272{,}220 & +0.2\% \\
    \bottomrule
  \end{tabular}
\end{table}

\subsection{Residual diagnostics}
\begin{figure}[H]
  \centering
  \includegraphics[width=0.95\linewidth]{04_residual_diagnostics.png}
  \caption{Holt--Winters in-sample residuals: time plot, histogram, and
  ACF.}
  \label{fig:fc-resid}
\end{figure}
The Ljung--Box statistic (\cref{eq:ljungbox}) at lag 12 is $Q_{LB}=16.11$,
$p=0.186$ — we fail to reject the null of no residual autocorrelation,
i.e.\ the fitted Holt--Winters model has not left obviously exploitable
structure in its residuals, a reasonable (if not definitive, given
$T=48$) check of adequacy.

\section{Interpretation}
Three points carry forward. First, the ADF results
(\cref{tab:adf}) justify the modelling choice made throughout this
chapter: the series' seasonal \emph{shape}, not only its trend, needs
explicit estimation, which is exactly what both Holt--Winters'
seasonal state $s_t$ and SARIMA's $(P,D,Q)_{12}$ terms provide, and is
why a naive linear-trend-only model is not attempted. Second, the
disagreement between the single hold-out split and the 8-fold rolling
backtest on which of Holt--Winters/SARIMA is "better" is not noise to be
explained away — it is the expected behaviour of a short series with
few independent forecast origins, and is the reason this report reports
\emph{both} evaluations rather than picking whichever looks best. Third,
the practical takeaway for the business is directional, not a claim of
precision: total sales are projected up by roughly a quarter next year,
but growth is sharply uneven by category — Office Supplies is
accelerating, Technology is flat — which is a materially different and
more actionable statement than a single store-wide growth number, and
motivates category-aware inventory/staffing planning rather than a
uniform +25\% assumption.
